% Part: applied-modal-logics
% Chapter: temporal-logic
% Section: possible-histories

\documentclass[../../../include/open-logic-section]{subfiles}

\begin{document}

\olfileid[hi]{aml}{tl}{poss}

\olsection{संभावित इतिहास}

अब तक प्रयुक्त कालिक तर्कशास्त्र के संबंधात्मक प्रतिरूप अत्यंत लचीले हैं,
क्योंकि अभिगम्यता संबंध पर कोई प्रतिबंध लगाना आवश्यक नहीं है। इसका अर्थ है
कि कालिक प्रतिरूप अतीत और भविष्य दोनों में शाखित हो सकते हैं; किंतु हम
शाखन की अधिक ``मोडल'' अवधारणा पर विचार करना चाह सकते हैं, जिसमें घटनाओं के
अनुक्रमों को संभावित इतिहास माना जाता है। इसके लिए हमारी भाषा बदलना अनिवार्य
नहीं है, यद्यपि हम अपने ``सामान्य'' मोडल संकारक $\Box$ और~$\Diamond$ भी
जोड़ सकते हैं। समय के साथ कर्ताओं के ज्ञान में होने वाले परिवर्तनों को
निरूपित करने के लिए हम ज्ञानात्मक अभिगम्यता संबंध जोड़ने पर भी विचार कर
सकते हैं।

\begin{defn}
  कालिक भाषा का \emph{संभावित-इतिहास प्रतिरूप} एक त्रिक
  $\mModel{M} = \tuple{T, C, V}$ है, जहाँ
  \begin{enumerate}
    \item $T$ एक अरिक्त समुच्चय है, जिसके अवयव कालिक अवस्थाएँ माने जाते हैं।
    \item $C$ किसी तंत्र के संगणनात्मक पथों, अथवा \emph{संभावित इतिहासों},
      का समुच्चय है। दूसरे शब्दों में, $C$ अवस्थाओं $s_1$, $s_2$,~$s_3$,
      \dots के अनुक्रमों~$\sigma$ का समुच्चय है, जहाँ प्रत्येक $s_i \in T$।
    \item $V$ ऐसा फलन है जो प्रत्येक !!{propositional variable}~$p$ को समय-बिंदुओं
      का समुच्चय~$V(p)$ देता है।
  \end{enumerate}
  बात को सरल रखने के लिए हम सामान्यतः यह भी मानेंगे कि जब कोई इतिहास
  $C$ में हो, तब उसके सभी उत्तरांश भी $C$ में हों। उदाहरण के लिए, यदि
  $s_1$, $s_2$,~$s_3$ कोई अनुक्रम $C$ में है, तो $s_2$, $s_3$ तथा केवल
  $s_3$ वाले अनुक्रम भी $C$ में हैं। साथ ही, जब दो अवस्थाएँ $s_i$ और~$s_j$
  किसी अनुक्रम~$\sigma$ में आती हैं, तब $i < j$ होने पर हम कहते हैं कि
  $s_i \prec_\sigma s_j$। जब $t \in V(p)$ हो, तब हम कहते हैं कि $p$,
  $t$ पर \emph{सत्य} है।
\end{defn}

एक प्रासंगिक परिवर्तन यह है कि प्रतिरूप~$\mModel M$ में किसी समय-बिंदु~$t$
पर किसी !!{formula} की सत्यता का मूल्यांकन अब हम ऐसे इतिहास~$\sigma$ के
सापेक्ष करते हैं जिसमें $t$ एक अवस्था के रूप में आता है। फिर भी
!!{propositional variable}ों या सत्यता-फलनीय संयोजकों के अर्थविज्ञान में कोई
परिवर्तन आवश्यक नहीं है। वे सभी ठीक वैसे ही हैं जैसे
\olref[sem]{defn:tmodels} में थे, क्योंकि उनमें से कोई भी~$\sigma$ का उल्लेख
नहीं करता। किंतु अब हम इन इतिहासों के सापेक्ष अपने भविष्य-संकारक~$\Ftemp$
को पुनः परिभाषित करते हैं और संकारक~$\Diamond$ जोड़ते हैं।

\begin{defn}\ollabel{defn:phmodels}
  $\mModel M$ में \emph{!!{formula}~$!A$ की $t, \sigma$ पर सत्यता}, जिसे
  $\mSat{M}{!A}[t, \sigma]$ से लिखते हैं:
  \begin{enumerate}
  \item\ollabel{defn:sub:phmodels-f} \indcase{!A}{\Ftemp
    !B}{$\mSat{M}{\indfrm}[t, \sigma]$ तभी जब $t \prec_\sigma t'$ वाले
    किसी $t' \in T$ के लिए $\mSat{M}{!B}[t', \sigma]$।}
  \item\ollabel{defn:sub:phmodels-diamond} \indcase{!A}{\Diamond
    !B}{$\mSat{M}{\indfrm}[t, \sigma]$ तभी जब किसी ऐसे $\sigma' \in C$
    के लिए $\mSat{M}{!B}[t, \sigma']$ हो जिसमें $t$ आता है।}
  \end{enumerate} 
\end{defn}

अन्य कालिक और मोडल संकारकों को भी इसी प्रकार परिभाषित किया जा सकता है। अब
हम ऐसे दावों को निरूपित कर सकते हैं जिनमें काल और मोडलता दोनों सम्मिलित हों।
उदाहरण के लिए, ``$p$ घटित नहीं होगा, किंतु वह घटित हो सकता था'' को हम
$\lnot \Ftemp p \land \Diamond \Ftemp p$ सूत्र से संकेतित कर सकते हैं। यह
सूत्र ऐसे बिंदु और इतिहास पर सत्य होगा जहाँ $p$ किसी उत्तरवर्ती अवस्था पर
सत्य नहीं होता, किंतु कोई वैकल्पिक इतिहास ऐसा है जहाँ $p$ भविष्य में सत्य
हो जाता है।

\end{document}
